nanaEEE nanaAAA


Una differenza fondamentale con~\cite{boucheron2005theory}, da cui nonostante questo prendiamo ispirazione per i risultati teorici, lies in the perspective from which bounds are analysed. Quel lavoro valuta i bound in base alla dipendenza da $n$; nei nostri esperimenti, invece, $n$ è sempre fissato e guardiamo ai bound come funzioni della complessità dello spazio delle ipotesi. Per esempio, risultati come il Teorema~\ref{thm: bound emprisk} che trova una formulazione analoga nel Teorema 5.1 di~\cite{boucheron2005theory}, a noi interessano perché ci permettono di includere il ruolo del rumore, mentre nel paper permettono di ottenere una dipendenza da $n$ migliore quando l'errore empirico va a $0$. Non solo, ma se non venisse combinato con il risultato di saturazione, questo bound ci darebbe una dipendenza dalla complessità del modello peggiore, che è lineare invece che sotto radice. Questa divergenza tra i lavori deriva dal fatto che loro cercano delle garanzie teoriche sulla precisione del rischio empirico nello stimare il rischio vero di un estimatore, le quali migliorano se si considera un dataset più grande, mentre per noi i bound sono funzioni che vogliamo descrivano bene l'andamento qualitativo della curva del gap tra rischio empirico e vero, all'aumentare della complessità del modello. 